社交平台上的信息扩散是显著\emph{重尾}的：极少数级联会触达巨大受众，而绝大
多数迅速消亡；并且经验研究已经确立，虚假内容比真实内容传播得更快、更深、
更广。本文追问由此产生的公众对误信息的易感性是否由一个\emph{共振}机制所
支配，并在一个空间耦合的双稳态意见模型中研究该问题，其随机驱动为
$\alpha$-稳定 \Levy\ 噪声。

我们结合三套互补的分析装置。(i) \emph{半解析的响应理论}：对空间分数阶
Fokker--Planck 方程做非微扰矩展开，得到作为幅值归一化共振指标的\emph{谱放大
因子}。(ii) \emph{逃逸率分析}：双稳态势的 Kramers 速率给出最优噪声水平的
无参数预言 $D^{*}=\Delta U/2$，我们用仿真并对照实测的平均首次穿越时间来
检验它。(iii) \emph{信息论分析}：输入--输出互信息结合禁区定理，确立了针对
不规则（而非仅周期）公共信号流的非周期随机共振。

在 $50\times50$ 格点（$N=2500$）上做系综平均的仿真中我们发现：高斯
（$\alpha=2$）与重尾（$\alpha=1.5$）噪声下均出现随机共振；重尾驱动下共振
显著\emph{右移并展宽}，我们将其追溯为跳跃介导的、而非热激活的势垒跨越；
存在一个\emph{误信息脆弱性窗口}，其中信噪比与集体信念切换的相干性在同一
噪声水平处同时达到峰值；以及辟谣具有强烈的非线性回报——真相偏置的边际效应
恰好在该窗口内最大。

我们认为这些结果带有明确的治理含义：一个高度同步、高度一致的公共领域并非
自明地健康，而共振框架同时识别出群体\emph{何时}最脆弱，以及干预\emph{何处}
最有效。

\medskip
\noindent\textbf{关键词：}误信息；意见动力学；社交网络；随机共振；重尾动力学；
\Levy\ 噪声；Kramers 逃逸；互信息；事实核查；级联动力学。

\medskip
\noindent\textbf{MSC 2020：}91D30, 60G51, 82C31, 60H10, 60J60.
